\documentclass[a4paper,12pt]{extarticle}
\usepackage[margin=1.5cm]{geometry}
\usepackage{amsmath, amssymb}
\usepackage{array}
\usepackage{multicol}
\usepackage{helvet}
\usepackage{multirow}
\usepackage{graphicx}
\renewcommand{\baselinestretch}{2}
\renewcommand{\familydefault}{\sfdefault}

\begin{document}

    \begin{tabular}{| c | c | c | c |}
        \hline
        & MATEMÁTICAS & 20/04/2026 & calificación \\
        \hline 

        \multirow{2}{*}{\includegraphics{../../img-Logos/iesoGalileo.jpg}} & Nombre y apellidos \hspace{120pt} & Ex. 2º ESO:  & \\ 
        &&Rec. 2º Evaluación& \\ 
        \hline
    \end{tabular}
    
    \begin{tabular}{| l | }
        \hline
        NOTAS: \\ 
            1. El examen tiene que ser limpio, ordenado y sin faltas de ortografía. \\
            
            2. El examen ha de realizarse en bolígrafo negro o azul, evitando tachones en la medida \\
            de lo posible. \\ 
            
            3. Deben aparecer todas las operaciones NO vale con indicar solo el resultado.\\
            
            4. Los problemas deben contener: Datos, planteamiento y resolución, respondiendo a lo \\
            que se pregunte, NO vale con indicar un número como solución del problema.\\

            5. Persona que se vea copiando o hablando con un compañero tiene un cero en el examen\\

            6. NO está permitido el uso de calculadoras \\
        \hline
    \end{tabular}

    \begin{enumerate}
        \item (1 punto) Completa la tabla con la siguiente información: 
        
        \begin{itemize}
            \item Cristina tiene x años
            \item Alberto, su esposo, tiene 3 años más 
            \item Javier, su padre, le dobla la edad
            \item Marta, su madre, tiene 4 años menos que su padre
            \item Loli y Mar son sus hijas gemelas. Las tuvo con 26 años.
            \item Javi, el pequeño, tiene la mitad de años que las gemelas.
        \end{itemize} 
        \begin{tabular}{| c | c | c | c | c | c | c | }
                \hline
                & Cristina & Alberto & Javier & Marta & Loli y Mar &  Javi \\ \hline
                Edad & x & $ \hspace{2cm}$  & $ \hspace{2cm}$ & $ \hspace{2cm}$  &$ \hspace{2cm}$ &  $ \hspace{2cm}$\\ 
                &  &  &  &   & &  \\ 
                \hline
            \end{tabular}
        \item (1 punto) Extrae factor común en cada una de las siguientes expresiones.
        \begin{enumerate}
            \item $8x^2+12x^3$
            \item $4x^2y+6xy^2+2xy$
        \end{enumerate}
        \item (1 punto) Desarrolla utilizando las identidades notables.
        \begin{enumerate}
            \item $(3x+5)^2$
            \item $(x-4)^2$
        \end{enumerate}
        \item (1 punto) Resuelve la ecuación de primer grado
        \[x-(x+4)-3(x+2)= 7-2(x+5)\]
        \item (1 punto) Resuelve la ecuación de segundo grado
        \[x^2+x-6=0\]
        \item (1 punto) Si al doble de un número se le resta su mitad, resulta 54. Hallar ese número.
        \item (0.5 puntos) Resuelve la ecuación de segundo grado incompleta
        \[9x^2-3x=0\]
        \item (0.5 puntos) Resuelve la ecuación de segundo grado incompleta
        \[2x^2-50=0\]
        \item (1,5 puntos) Resuelve el siguiente sistema
        $\left \{\begin{array}{rr}
                    x+3y=2 \\
                    -2x+3y=5
                \end{array}\right .$ 
        \item (1.5 puntos) En un patio de una escuela infantil hay triciclos y cochecitos. En total son 20 y todas las ruedas suman 67. ¿Cuántos triciclos y cuántos cochecitos hay en el patio?$\vspace{7cm}$ \newpage

    \end{enumerate}
\end{document}